\documentclass[12pt]{article}

\usepackage{sbc-template}
\usepackage{graphicx}
\usepackage{url}
\usepackage[brazil]{babel}
\usepackage[utf8]{inputenc}
\usepackage{booktabs}
\usepackage[pdftex]{hyperref}
\usepackage{indentfirst}
\usepackage{enumitem}

\sloppy

% ============================================================================
% MODELO DIDÁTICO DE ARTIGO
% Sistemas de Apoio à Decisão (SAD) + Design Science Research (DSR)
% + Design Centrado no Usuário (DCU) + Engenharia de Software
% + UML / Modelo de Visões 4+1 + Avaliação ISO/IEC 25010 + SUS + TAM
%
% Orientação geral:
% - Substitua os textos de orientação pelo conteúdo efetivo do trabalho.
% - Mantenha evidências para cada decisão de projeto.
% - O artigo deve mostrar a rastreabilidade:
%   PROBLEMA -> USUÁRIO -> REQUISITOS -> ARQUITETURA -> ARTEFATO -> AVALIAÇÃO.
% ============================================================================

\title{Título do Sistema de Apoio à Decisão: Subtítulo}

\author{Nome do Autor Um\inst{1}, Nome do Autor Dois\inst{1}, \\
Nome do Autor Três\inst{1}, Nome do Autor Quatro\inst{1}, \\
Nome do Autor N.\inst{1}, Nome do Coorientador\inst{2},
Gilton José Ferreira da Silva\inst{2}}

\address{Departamento de Computação (DCOMP)\\
Universidade Federal de Sergipe (UFS)\\
Av. Marechal Rondon, s/n -- Jardim Rosa Elze -- CEP 49100-000\\
São Cristóvão -- SE -- Brazil
\nextinstitute
Programa de Pós-Graduação em Ciência da Computação (PROCC)\\
Universidade Federal de Sergipe (UFS)\\
Av. Marechal Rondon, s/n -- Jardim Rosa Elze -- CEP 49100-000\\
São Cristóvão -- SE -- Brazil
\email{autor01@dcomp.ufs.br, autor02@dcomp.ufs.br}
\email{autor03@dcomp.ufs.br, autor04@dcomp.ufs.br}
\email{gilton@dcomp.ufs.br}
}

\begin{document}

\maketitle

\begin{abstract}
\textbf{Context:} Briefly present the decision problem, application domain and target decision-maker.
\textbf{Objective:} State the objective of designing, developing and evaluating the Decision Support System (DSS).
\textbf{Method:} Report Design Science Research as the methodological framework, complemented by User-Centered Design, Software Engineering, UML and the 4+1 View Model. Mention the main evaluation instruments adopted.
\textbf{Results:} Summarize the developed artifact and the principal technical, usability and technology-acceptance results.
\textbf{Conclusions:} State the main practical and scientific contributions and the principal limitations.
\end{abstract}

\begin{resumo}
\textbf{Contexto:} Apresente, em poucas frases, o domínio, o problema decisório, o tomador de decisão e por que existe necessidade de apoio computacional.
\textbf{Objetivo:} Declare o objetivo geral, preferencialmente com os verbos ``projetar, desenvolver e avaliar'' um Sistema de Apoio à Decisão (SAD).
\textbf{Método:} Informe a adoção da Design Science Research (DSR), integrada ao Design Centrado no Usuário (DCU), Engenharia de Software, UML e Modelo de Visões 4+1. Cite resumidamente como o artefato foi avaliado.
\textbf{Resultados:} Sintetize o artefato construído e os principais resultados da avaliação técnica, SUS e TAM.
\textbf{Conclusões:} Apresente a contribuição do SAD para o problema decisório e as principais limitações.
\end{resumo}

% ============================================================================
\section{Introdução}
\label{sec:introducao}

\textbf{O que escrever nesta seção:}
Apresente o domínio de aplicação e conduza o leitor do contexto geral até o problema específico de decisão. Caracterize quem decide, quais decisões precisam ser tomadas, quais dados estão disponíveis, quais dificuldades existem no processo atual e quais consequências decorrem dessas dificuldades.

Explique por que um Sistema de Apoio à Decisão é adequado ao problema. Evite iniciar pela tecnologia. O foco inicial deve estar no \textbf{problema decisório}. Sustente a argumentação com literatura científica, dados, documentos institucionais ou outras evidências.

Ao final da Introdução, o leitor deve compreender claramente:
(i) qual é o problema;
(ii) quem é afetado;
(iii) por que o problema é relevante;
(iv) por que um SAD pode contribuir; e
(v) qual lacuna o trabalho pretende tratar.

\subsection{Problema de Pesquisa e Questão de Pesquisa}

Formule explicitamente o problema investigado. Em DSR, é recomendável que a questão esteja relacionada à criação e avaliação de um artefato.

\textbf{Exemplo:}
``Como projetar, desenvolver e avaliar um Sistema de Apoio à Decisão capaz de apoiar [tipo de usuário] na decisão sobre [decisão], considerando [restrições/variáveis relevantes]?''

\subsection{Justificativa}

Demonstre a relevância prática e científica da pesquisa. Utilize evidências que indiquem que o problema existe e que merece intervenção. Diferencie, quando possível, a relevância para o usuário/organização da contribuição esperada para a área de Sistemas de Apoio à Decisão.

\subsection{Objetivos}

Apresente um objetivo geral e objetivos específicos.

\textbf{Exemplo de objetivo geral:}
Projetar, desenvolver e avaliar um Sistema de Apoio à Decisão para apoiar [tomador de decisão] em [decisão], utilizando Design Science Research como abordagem metodológica.

\textbf{Exemplos de objetivos específicos:}
\begin{itemize}[noitemsep]
    \item compreender o problema decisório e as necessidades dos usuários;
    \item identificar requisitos funcionais, não funcionais e requisitos de decisão;
    \item projetar o SAD utilizando UML e o Modelo de Visões 4+1;
    \item desenvolver um artefato funcional;
    \item avaliar tecnicamente a solução;
    \item avaliar sua usabilidade por meio do System Usability Scale (SUS);
    \item analisar sua aceitação por meio do Technology Acceptance Model (TAM).
\end{itemize}

\subsection{Organização do Artigo}

Descreva brevemente o conteúdo das seções seguintes. Use apenas um parágrafo.

% ============================================================================
\section{Fundamentação Teórica}
\label{sec:fundamentacao}

Esta seção deve apresentar somente os conceitos necessários para compreender o problema, a construção do artefato e sua avaliação. Evite transformar a fundamentação em um glossário.

\subsection{Sistemas de Apoio à Decisão}

Fundamente o conceito de SAD a partir da literatura da área, incluindo Turban e colaboradores. Caracterize o tipo de decisão tratada (estruturada, semiestruturada ou não estruturada), os usuários/tomadores de decisão e os componentes relevantes do SAD.

Quando aplicável, descreva:
\begin{itemize}[noitemsep]
    \item subsistema de gerenciamento de dados;
    \item subsistema de gerenciamento de modelos;
    \item interface/interação com o usuário;
    \item conhecimento/regras de negócio;
    \item processo decisório apoiado pelo sistema.
\end{itemize}

Explique explicitamente \textbf{por que o sistema desenvolvido pode ser classificado como um SAD}, e não apenas como um sistema de informação convencional.

\subsection{Design Science Research}

Apresente a DSR como abordagem voltada à criação e avaliação de artefatos destinados a solucionar problemas relevantes. Fundamente a abordagem adotada e informe qual processo/ciclo de DSR será operacionalizado no trabalho.

Explique os conceitos de problema, artefato, rigor, relevância, construção e avaliação. Esta fundamentação deve preparar o leitor para a Seção~\ref{sec:metodologia}.

\subsection{Design Centrado no Usuário}

Apresente os princípios de Design Centrado no Usuário utilizados no trabalho. Explique a função de técnicas como identificação de stakeholders, persona, mapa de empatia e, quando pertinente, jornada do usuário.

Deixe claro que essas técnicas serão utilizadas para compreender o tomador de decisão e alimentar a elicitação de requisitos.

\subsection{Engenharia de Software, UML e Modelo de Visões 4+1}

Apresente resumidamente a Engenharia de Software e a UML como suporte à especificação e ao projeto do artefato.

Explique o Modelo de Visões 4+1 de Philippe Kruchten:
\begin{itemize}[noitemsep]
    \item \textbf{Visão de Casos de Uso (+1):} cenários e funcionalidades observáveis pelos atores;
    \item \textbf{Visão Lógica:} estrutura e responsabilidades do domínio;
    \item \textbf{Visão de Processo:} comportamento, fluxos, concorrência e comunicação;
    \item \textbf{Visão de Desenvolvimento:} organização dos módulos/componentes de software;
    \item \textbf{Visão Física:} implantação do sistema em infraestrutura computacional.
\end{itemize}

Informe que as visões serão representadas por diagramas UML adequados ao contexto do projeto.

\subsection{Avaliação de Sistemas}

Apresente os fundamentos dos instrumentos que serão utilizados posteriormente:
qualidade técnica de software, SUS e TAM. Não apresente ainda os resultados.

% ============================================================================
\section{Trabalhos Relacionados}
\label{sec:relacionados}

Realize uma análise crítica de trabalhos que desenvolveram soluções semelhantes para o mesmo domínio ou problema decisório.

Para cada trabalho, procure identificar: problema, público, tipo de SAD, dados utilizados, método/modelo decisório, tecnologia, estratégia de avaliação e resultados.

Ao final, apresente uma síntese comparativa, preferencialmente em tabela, e explicite a \textbf{lacuna} que diferencia ou justifica o artefato proposto.

% ============================================================================
\section{Metodologia}
\label{sec:metodologia}

Este trabalho adota a \textit{Design Science Research} como eixo metodológico. O processo integra técnicas de Design Centrado no Usuário e Engenharia de Software para construir e avaliar um Sistema de Apoio à Decisão.

\textbf{Importante:} esta seção descreve \textbf{como a pesquisa foi conduzida}. Os resultados produzidos por cada etapa devem ser apresentados posteriormente.

\subsection{Visão Geral do Processo de Design Science Research}

Apresente uma figura com o fluxo metodológico adotado. O modelo didático deste trabalho pode ser organizado nas seguintes etapas:

\begin{enumerate}
    \item identificação e compreensão do problema;
    \item definição dos objetivos da solução;
    \item projeto e desenvolvimento do artefato;
    \item demonstração do SAD;
    \item avaliação do artefato;
    \item comunicação dos resultados.
\end{enumerate}

Explique que o processo pode ser iterativo: resultados de avaliação podem conduzir a revisões no problema, requisitos, projeto ou implementação.

\subsection{Etapa 1 -- Identificação e Compreensão do Problema}

Descreva como o problema decisório foi investigado: revisão de literatura, documentos, entrevistas, observação, questionários, dados organizacionais ou outras fontes.

\subsubsection{Stakeholders}

Identifique os atores relacionados ao processo decisório, destacando o \textbf{tomador de decisão principal}, usuários secundários e demais interessados.

\subsubsection{Persona}

Construa uma persona representativa do usuário/tomador de decisão. Inclua objetivos, contexto, responsabilidades, conhecimentos, necessidades, dificuldades e comportamento relevante para o uso do SAD.

\subsubsection{Mapa de Empatia do Tomador de Decisão}

Apresente o mapa de empatia e explique, no texto, as principais evidências levantadas: o que o usuário pensa e sente, vê, ouve, fala e faz, além de suas dores e ganhos esperados.

Não utilize o mapa apenas como elemento visual: demonstre como suas descobertas influenciaram os requisitos do sistema.

\subsubsection{Processo Decisório Atual}

Descreva como a decisão é tomada atualmente (\textit{AS-IS}). Identifique entradas, atores, critérios, etapas, dificuldades, gargalos e saídas. Quando pertinente, represente o processo por diagrama de atividades, BPMN ou outro modelo adequado.

\subsection{Etapa 2 -- Definição dos Objetivos da Solução}

Converta os problemas identificados em objetivos mensuráveis para o artefato.

\textbf{Exemplo:} reduzir o tempo necessário para analisar alternativas; aumentar a consistência da decisão; consolidar informações dispersas; permitir comparação multicritério; fornecer explicações para recomendações.

Defina critérios que permitam avaliar posteriormente se esses objetivos foram alcançados.

\subsection{Etapa 3 -- Projeto e Desenvolvimento}

Descreva o processo de Engenharia de Software adotado, as tecnologias e a estratégia de desenvolvimento. Explique como ocorreu a passagem das necessidades dos usuários para requisitos, arquitetura, protótipo e implementação.

\subsubsection{Elicitação e Especificação de Requisitos}

Apresente os métodos utilizados para elicitar e validar requisitos.

Organize, preferencialmente:
\begin{itemize}[noitemsep]
    \item requisitos funcionais;
    \item requisitos não funcionais;
    \item regras de negócio;
    \item requisitos relacionados ao processo de decisão;
    \item fontes e requisitos de dados.
\end{itemize}

Utilize identificadores (RF01, RF02, RNF01 etc.) para permitir rastreabilidade.

\subsubsection{Prototipação e Design da Interface}

Explique como foram produzidos os protótipos e como as necessidades identificadas em persona/mapa de empatia foram refletidas nas interfaces.

\subsubsection{Modelagem UML e Arquitetura 4+1}

Explique quais diagramas serão utilizados em cada visão e justifique sua seleção. Não é obrigatório utilizar todos os diagramas UML existentes.

Sugestão didática:
\begin{itemize}[noitemsep]
    \item \textbf{Casos de Uso (+1):} diagrama de casos de uso e descrição dos cenários principais;
    \item \textbf{Visão Lógica:} diagrama de classes e, quando necessário, diagramas de sequência;
    \item \textbf{Visão de Processo:} diagramas de atividades e/ou sequência para os fluxos críticos;
    \item \textbf{Visão de Desenvolvimento:} diagrama de componentes e/ou pacotes;
    \item \textbf{Visão Física:} diagrama de implantação.
\end{itemize}

\subsubsection{Mecanismo de Apoio à Decisão}

Esta subseção é essencial para caracterizar o artefato como SAD.

Descreva como o sistema transforma dados em apoio à decisão: modelos analíticos, regras, indicadores, métodos multicritério, simulação, otimização, aprendizado de máquina, análise ``what-if'', recomendação ou outra técnica.

Apresente entradas, processamento, critérios/pesos, saídas e a forma como o resultado é oferecido ao tomador de decisão.

\subsection{Etapa 4 -- Demonstração}

Descreva o cenário no qual o artefato será demonstrado. Defina tarefas representativas e casos de uso que permitam observar o SAD apoiando uma decisão realista.

\subsection{Etapa 5 -- Plano de Avaliação do Artefato}

A avaliação deve responder se o artefato foi construído adequadamente e se produz utilidade no contexto para o qual foi projetado. Recomenda-se combinar diferentes dimensões de avaliação.

\subsubsection{Avaliação Técnica -- Qualidade de Software}

Utilize a ISO/IEC 25010 como referência para estruturar a avaliação técnica. Selecione somente características pertinentes ao artefato e explique como serão operacionalizadas.

Exemplos de características:
\begin{itemize}[noitemsep]
    \item adequação funcional;
    \item eficiência de desempenho;
    \item confiabilidade;
    \item segurança;
    \item compatibilidade;
    \item manutenibilidade;
    \item portabilidade.
\end{itemize}

Para cada característica escolhida, defina uma métrica ou procedimento verificável.

\begin{table}[!ht]
\centering
\caption{Exemplo de plano de avaliação técnica}
\label{tab:avaliacao_tecnica}
\begin{tabular}{p{3.0cm}p{5.0cm}p{3.0cm}}
\toprule
\textbf{Característica} & \textbf{Métrica/Procedimento} & \textbf{Critério} \\
\midrule
Adequação funcional & Percentual de requisitos funcionais aprovados nos testes & Definir meta \\
Desempenho & Tempo de resposta em tarefas críticas & Definir limite \\
Confiabilidade & Taxa de execução sem falhas nos cenários de teste & Definir meta \\
Segurança & Verificação dos controles pertinentes ao contexto & Definir critério \\
\bottomrule
\end{tabular}
\end{table}

Quando o SAD possuir modelo preditivo, recomendador ou algoritmo específico, inclua também métricas próprias da técnica utilizada. A ISO/IEC 25010 não substitui a validação do modelo decisório.

\subsubsection{Avaliação da Qualidade do Apoio à Decisão}

Defina indicadores relacionados à finalidade do SAD, tais como:
qualidade/consistência da recomendação, tempo para tomada de decisão, redução de esforço, comparação de alternativas, disponibilidade das informações necessárias e aderência às regras do domínio.

Quando houver um ``resultado correto'' conhecido, utilize métricas adequadas. Quando a decisão depender de julgamento humano, considere cenários, especialistas ou comparação com o processo anterior.

\subsubsection{Avaliação de Usabilidade -- System Usability Scale (SUS)}

A usabilidade será avaliada por meio do \textit{System Usability Scale} (SUS), instrumento proposto por John Brooke em 1986.

O SUS possui dez afirmações respondidas em escala Likert de cinco pontos, de forte discordância a forte concordância. Recomenda-se aplicar o instrumento após os participantes utilizarem o SAD em tarefas representativas.

\textbf{Exemplo de instrumento SUS:}
\begin{enumerate}
    \item Eu acho que gostaria de usar este sistema com frequência.
    \item Eu achei o sistema desnecessariamente complexo.
    \item Eu achei o sistema fácil de usar.
    \item Eu acho que precisaria do apoio de uma pessoa com conhecimentos técnicos para usar este sistema.
    \item Eu achei que as várias funções deste sistema estavam bem integradas.
    \item Eu achei que havia muita inconsistência neste sistema.
    \item Eu imagino que a maioria das pessoas aprenderia a usar este sistema rapidamente.
    \item Eu achei o sistema muito complicado de usar.
    \item Eu me senti confiante ao usar o sistema.
    \item Eu precisei aprender muitas coisas antes de conseguir utilizar o sistema.
\end{enumerate}

\textbf{Cálculo:}
para os itens ímpares, subtraia 1 da resposta; para os itens pares, subtraia a resposta de 5. Some as dez contribuições e multiplique o resultado por 2,5. O escore final varia de 0 a 100.

Apresente o escore global e, quando pertinente, estatísticas descritivas. Não trate o valor do SUS como uma porcentagem de usabilidade.

\subsubsection{Avaliação de Aceitação -- Technology Acceptance Model (TAM)}

Utilize o TAM para investigar a aceitação do SAD. Como estrutura mínima para a disciplina, avalie:
\textbf{Utilidade Percebida (UP)}, \textbf{Facilidade de Uso Percebida (FUP)} e \textbf{Intenção de Uso (IU)}.

As afirmações abaixo são exemplos didáticos e devem ser adaptadas ao contexto do SAD e fundamentadas no instrumento/literatura adotados.

\begin{table}[!ht]
\centering
\caption{Exemplo de instrumento baseado no TAM}
\label{tab:tam}
\begin{tabular}{p{2.0cm}p{9.2cm}}
\toprule
\textbf{Construto} & \textbf{Exemplo de afirmação} \\
\midrule
UP1 & Utilizar o SAD melhora minha capacidade de tomar a decisão analisada. \\
UP2 & O SAD torna o processo de decisão mais eficiente. \\
UP3 & O SAD é útil para as atividades relacionadas a essa decisão. \\
FUP1 & Aprender a utilizar o SAD é fácil para mim. \\
FUP2 & Minha interação com o SAD é clara e compreensível. \\
FUP3 & Considero o SAD fácil de utilizar. \\
IU1 & Eu pretendo utilizar este SAD quando precisar tomar esse tipo de decisão. \\
IU2 & Se o SAD estiver disponível, pretendo incorporá-lo ao meu processo de trabalho. \\
IU3 & Eu recomendaria a utilização do SAD em situações semelhantes. \\
\bottomrule
\end{tabular}
\end{table}

Utilize uma escala Likert definida previamente (por exemplo, cinco ou sete pontos) e mantenha a mesma escala para todos os itens. Explique como os escores dos construtos serão calculados.

Para análises exploratórias em turmas com amostras pequenas, podem ser apresentadas médias/medianas e dispersão por construto. Se o trabalho pretender testar formalmente relações causais do TAM, deverá utilizar instrumento validado, amostra adequada e técnica estatística compatível.

\subsubsection{Avaliação por Especialistas (Opcional)}

Quando pertinente, especialistas no domínio podem avaliar a correção das regras, indicadores, recomendações ou resultados do SAD. Informe critérios de seleção dos especialistas e o protocolo utilizado.

\subsubsection{Síntese e Triangulação das Evidências}

Integre as diferentes avaliações:
\begin{itemize}[noitemsep]
    \item \textbf{ISO/IEC 25010 e testes:} o software atende aos critérios técnicos?
    \item \textbf{Qualidade do apoio à decisão:} o artefato melhora ou apoia efetivamente a decisão?
    \item \textbf{SUS:} o sistema apresenta usabilidade adequada?
    \item \textbf{TAM:} os usuários percebem utilidade, facilidade e intenção de uso?
    \item \textbf{Especialistas:} o conteúdo/modelo decisório é coerente com o domínio?
\end{itemize}

A triangulação deve permitir discutir não apenas se ``o sistema funciona'', mas se o artefato é tecnicamente adequado, utilizável, aceito e relevante para o problema.

\subsection{Etapa 6 -- Comunicação}

Informe que os resultados do ciclo de DSR serão consolidados e comunicados no próprio artigo, incluindo descrição do artefato, evidências de avaliação, limitações e contribuição.

% ============================================================================
\section{Desenvolvimento e Descrição do Artefato}
\label{sec:artefato}

Nesta seção, apresente os \textbf{resultados da construção}. Diferentemente da metodologia, aqui deve ser mostrado o que efetivamente foi produzido.

\subsection{Visão Geral do SAD}

Apresente nome, objetivo, usuários, contexto de uso, entradas, principais funcionalidades e saídas do sistema.

Inclua uma figura de visão geral, quando possível.

\subsection{Requisitos Obtidos}

Apresente os principais requisitos e demonstre sua relação com as necessidades identificadas no estudo com usuários.

Uma matriz de rastreabilidade pode relacionar:
\textit{dor/necessidade -> requisito -> funcionalidade -> teste/avaliação}.

\subsection{Arquitetura e Modelagem 4+1}

Apresente os diagramas produzidos para cada visão e explique as principais decisões arquiteturais.

\subsubsection{Visão de Casos de Uso (+1)}
Apresente atores, casos de uso e cenários principais.

\subsubsection{Visão Lógica}
Apresente a organização conceitual do sistema, classes/entidades e relações relevantes.

\subsubsection{Visão de Processo}
Apresente os fluxos dinâmicos críticos, especialmente aqueles relacionados à geração do apoio à decisão.

\subsubsection{Visão de Desenvolvimento}
Apresente módulos, componentes, camadas, serviços e dependências.

\subsubsection{Visão Física}
Apresente a implantação do sistema, servidores, banco de dados, dispositivos e integrações relevantes.

\subsection{Implementação do Mecanismo de Apoio à Decisão}

Apresente detalhadamente o mecanismo decisório implementado e demonstre, por exemplo, como uma entrada é transformada em recomendação, classificação, ranking, indicador, simulação ou análise.

\subsection{Interfaces e Fluxo de Uso}

Apresente as principais telas em ordem coerente com a jornada do usuário e com os casos de uso. Explique como as decisões de interface respondem às necessidades identificadas na etapa de DCU.

% Exemplo:
% \begin{figure}[!ht]
% \centering
% \includegraphics[width=.9\linewidth]{images/telas.jpg}
% \caption{Principais interfaces do Sistema de Apoio à Decisão.}
% \label{fig:telas}
% \end{figure}

% ============================================================================
\section{Avaliação e Resultados}
\label{sec:resultados}

Apresente os resultados conforme o plano definido na metodologia. Informe participantes, ambiente, tarefas, dados coletados e condições de execução antes de interpretar os resultados.

\subsection{Resultados da Avaliação Técnica}

Apresente os resultados das métricas selecionadas da ISO/IEC 25010, testes funcionais e métricas específicas do mecanismo decisório.

Não escreva apenas ``o sistema funcionou''. Mostre evidências quantitativas ou qualitativas verificáveis.

\subsection{Resultados da Avaliação do Apoio à Decisão}

Demonstre se o SAD atingiu os objetivos definidos para a decisão. Compare, quando possível, o cenário com e sem o artefato ou utilize cenários de referência/especialistas.

\subsection{Resultados do SUS}

Apresente número de participantes, escore SUS individual ou agregado, média/mediana, dispersão e interpretação fundamentada na literatura adotada.

\subsection{Resultados do TAM}

Apresente os resultados por construto (UP, FUP e IU). Utilize tabelas ou gráficos para facilitar a interpretação. Discuta quais dimensões apresentaram maior e menor concordância.

\subsection{Triangulação e Discussão}

Integre os resultados técnicos, decisórios, de usabilidade e aceitação.

Discuta:
\begin{itemize}[noitemsep]
    \item quais objetivos da solução foram atingidos;
    \item quais evidências sustentam a utilidade do artefato;
    \item possíveis contradições entre as avaliações;
    \item comparação com trabalhos relacionados;
    \item implicações para o domínio e para a área de SAD.
\end{itemize}

% ============================================================================
\section{Ameaças à Validade e Limitações}
\label{sec:limitacoes}

Discuta fatores que limitam a interpretação e generalização dos resultados.

Considere, quando pertinente:
\begin{itemize}[noitemsep]
    \item tamanho e perfil da amostra;
    \item seleção dos participantes;
    \item ambiente de avaliação;
    \item uso de dados simulados ou históricos;
    \item maturidade do protótipo;
    \item limitações das métricas;
    \item viés dos pesquisadores ou especialistas;
    \item adaptação de instrumentos;
    \item limitações do modelo decisório.
\end{itemize}

Diferencie limitações do artefato de ameaças à validade da avaliação.

% ============================================================================
\section{Considerações Finais}
\label{sec:consideracoes}

Retome o problema e o objetivo geral e sintetize como o trabalho os abordou.

Apresente as principais contribuições:
\begin{itemize}[noitemsep]
    \item \textbf{artefato:} SAD desenvolvido;
    \item \textbf{prática:} contribuição para o processo decisório;
    \item \textbf{conhecimento:} aprendizados, princípios, modelo ou evidências produzidas pela pesquisa.
\end{itemize}

Resuma os principais resultados da avaliação sem repetir toda a seção anterior. Informe as limitações mais relevantes e proponha trabalhos futuros coerentes com os achados.

% ============================================================================
\section*{Agradecimentos}
\label{sec:agradecimentos}

Agradeça às pessoas e instituições que contribuíram com a pesquisa, mas não se qualificam para autoria. Informe financiamentos e apoios institucionais quando exigido pelo evento ou periódico.

% ============================================================================
% REFERÊNCIAS SUGERIDAS PARA O ARQUIVO referencias.bib
%
% O estudante deve consultar as obras originais e preencher os dados
% bibliográficos completos no arquivo .bib.
%
% Sugestões de autores/obras a fundamentar:
% - Turban et al. -- Decision Support / Business Intelligence / Analytics.
% - Hevner et al. -- Design Science in Information Systems Research.
% - Peffers et al. -- Design Science Research Methodology (DSRM).
% - Kruchten -- Architectural Blueprints: The 4+1 View Model of Software Architecture.
% - Brooke -- System Usability Scale (SUS), 1986.
% - Davis -- Technology Acceptance Model (TAM).
% - ISO/IEC 25010 -- Systems and software Quality Requirements and Evaluation (SQuaRE).
% - Referências de Design Centrado no Usuário e UML adotadas na disciplina.
% ============================================================================

\bibliographystyle{sbc}
\bibliography{referencias}

\end{document}
